O Quadro~\ref{quad:comentarios_integrais} transcreve integralmente as respostas a A01
e A02. São quatro respostas em A01 e três em A02, dadas por quatro pessoas. O código
R corresponde à ordem do registro na exportação (1 a 12) e não identifica o
participante. Nenhuma resposta continha nome, e-mail ou outro dado pessoal, e não
houve supressões. A redação foi preservada, incluindo grafia e pontuação; a
marcação [\textit{sic}] indica grafia original.

\begin{quadro}[htb]
\centering
\legenda[quad:comentarios_integrais]{Transcrição integral dos comentários A01 e A02}
\footnotesize
\begin{tabular}{|P{1.1cm}|P{1.0cm}|P{10.6cm}|}
\hline
Código & Campo & Texto integral \\
\hline\hline
R03 & A01 & Fácil - explica bem o que devemos fazer, juntamente ao ensinamento que
repassa // Dificuldade - por mais que o fluxo esteja bem definido, as vezes causa
confusão para executar alguma coisa ou sair de um certo lugar. Mas pode ser
considerado dificuldade do próprio usuário em particular \\
\hline
R03 & A02 & Não se encaixa como limitação, mas pelo fato de ser uma solução para
aprender/auxiliar na programação, é válido ter algo muito bem explicado para
garantir a compreensão. Textos curtos as vezes não são adequados para ensinar, ao
mesmo tempo que os longos também não, mas se fosse algo mais didático seria o ideal.
Não vejo isso como limitação, mas sim como algo a ser levado em conta para ser
implementado se fizer sentido. Hoje a solução se encontra em um ótimo estado, mas a
aplicação depende do intuito e possibilidades do desenvolvedor do aplicativo \\
\hline
R05 & A01 & Muito bom, não tive dúvidas de utilização. \\
\hline
R05 & A02 & Uma sugestão de melhoria: O teste de dividir por zero não sinalizou que o
erro era a divisão por zero e apenas sinalizou a linha com a lógica errada. \\
\hline
R09 & A01 & Foi fácil criar um arquivo e começar a programar, bem simples. Foi
difícil achar o botão de debutar [\textit{sic}] (não entendi o que é trace) \\
\hline
R12 & A01 & Depuração de código. // Digitações na tela pequena. \\
\hline
R12 & A02 & Arrumar um jeito melhor de poder digitar num celular de forma mais fácil. \\
\hline
\end{tabular}
\vspace{0.2cm}
\fonte{Dados da pesquisa (2026). Barras duplas (//) indicam quebra de linha na resposta original.}
\end{quadro}
\FloatBarrier
